\section{交换环上的多项式环}

本节讨论交换幺环。若$R$是$R'$的子环，称$R'$为$R$的\textbf{扩环}；本节涉及扩环时，总假定$R,R'$有相同的单位元。

\begin{definition}
设$R'$为$R$的扩环，$u\in R'$。所有形如
\[
a_0+a_1u+\cdots+a_nu^n,\qquad a_i\in R,\quad n\ge0,
\]
的元素组成的集合记为$R[u]$，称为$u$在$R$上生成的子环，其中$u^0=1$。
\end{definition}
\begin{proof}[验证]
$R[u]$包含$0$，且对减法封闭。由交换性，
\begin{align*}
\left(\sum_{i=0}^n a_i u^i\right)
\left(\sum_{j=0}^m b_j u^j\right)
&=\sum_{i=0}^n\sum_{j=0}^m a_i b_j u^{i+j}\\
&=\sum_{k=0}^{n+m}c_k u^k,
\qquad c_k=\sum_{i+j=k}a_i b_j\in R.
\end{align*}
因此$R[u]$也对乘法封闭，是$R'$的子环。取常数多项式可知$R\subseteq R[u]$，取一次式$u$可知$u\in R[u]$。反过来，包含$R$与$u$的任何子环均包含$u$的所有非负整数次幂及它们的有限线性组合，故均包含$R[u]$。所以$R[u]$是包含$R$和$u$的最小子环。
\end{proof}

$R[u]$中的两个多项式表达式可能表示同一元素，这取决于$u$在$R$上满足的代数关系。先构造一个扩环及其中的元素$x$，使$x$的多项式只有在对应系数全相等时才相等。

\begin{definition}
用$R[[x]]$表示所有无穷序列
\[
(a_n)=(a_0,a_1,a_2,\ldots),\qquad a_i\in R,
\]
组成的集合，定义运算
\[
(a_n)+(b_n)=(a_n+b_n),\qquad (a_n)(b_n)=(c_n),
\]
其中
\[
c_n=a_0b_n+a_1b_{n-1}+\cdots+a_nb_0
=\sum_{i+j=n}a_i b_j.
\]
所得环称为$R$上的\textbf{一元形式幂级数环}，其元素称为形式幂级数。这里只使用序列及有限的系数计算，不涉及无穷级数的收敛性。
\end{definition}
% 原文将环公理的验证留作习题，以下补全。
\begin{proof}[验证]
逐项加法构成交换群，零元为$(0,0,\ldots)$，负元逐项取负。乘法的交换律由$a_i b_j=b_j a_i$及求和指标的互换得到。

设三个序列的系数分别为$a_i,b_j,c_k$，则两种结合次序所得第$n$项分别为
\[
\sum_{r+k=n}\left(\sum_{i+j=r}a_i b_j\right)c_k,
\qquad
\sum_{i+t=n}a_i\left(\sum_{j+k=t}b_j c_k\right).
\]
它们都等于有限和$\sum_{i+j+k=n}a_i b_j c_k$，所以乘法结合。又有
\[
\sum_{i+j=n}a_i(b_j+c_j)
=\sum_{i+j=n}a_i b_j+\sum_{i+j=n}a_i c_j,
\]
故分配律成立。单位元为$(1,0,0,\ldots)$，所以$R[[x]]$为交换幺环。

它包含子环$R_0=\{(a,0,0,\ldots):a\in R\}$，且$a\mapsto(a,0,0,\ldots)$给出$R\cong R_0$。将两者等同后，可把$R$看作$R[[x]]$的子环，此时
\[
a(b_0,b_1,\ldots)=(ab_0,ab_1,\ldots).
\]
\end{proof}

\begin{definition}
在$R[[x]]$中取$x=(0,1,0,\ldots)$，$x$在$R$上生成的子环$R[x]$称为$R$上的\textbf{一元多项式环}。
\end{definition}

按序列乘法，$x^i$只有第$i$个分量为$1$，其余分量均为$0$，这里从第$0$个分量开始计数。因此
\[
a_0+a_1x+\cdots+a_nx^n=(a_0,a_1,\ldots,a_n,0,\ldots).
\]
所以$R[x]$恰由只有有限个非零分量的序列组成，而且
\[
a_0+a_1x+\cdots+a_nx^n=0
\quad\Longleftrightarrow\quad a_0=\cdots=a_n=0.
\]
两个多项式相等当且仅当对应系数相等；比较时允许在末尾添上零系数。称$x$为$R$上的一个\textbf{未定元}，也称不定元。

\begin{theorem}\label{p2:poly-universal}
设$\sigma:R\to S$为交换幺环之间的环同态，且$\sigma(1_R)=1_S$。对任意$u\in S$，$\sigma$可以唯一地扩充为环同态
\[
\sigma_u:R[x]\to S,
\]
使$\sigma_u(x)=u$。具体地，
\[
\sigma_u\left(\sum_{i=0}^n a_i x^i\right)
=\sum_{i=0}^n\sigma(a_i)u^i.
\]
\end{theorem}
\begin{proof}
由系数表达式的唯一性，上式给出一个良好定义的映射。设$f(x)=\sum_i a_i x^i$、$g(x)=\sum_i b_i x^i$，补齐零系数后，
\begin{align*}
\sigma_u(f+g)
&=\sum_i\sigma(a_i+b_i)u^i\\
&=\sum_i\sigma(a_i)u^i+\sum_i\sigma(b_i)u^i\\
&=\sigma_u(f)+\sigma_u(g).
\end{align*}
再设$fg=\sum_k c_kx^k$，其中$c_k=\sum_{i+j=k}a_i b_j$。则
\begin{align*}
\sigma_u(fg)
&=\sum_k\sigma(c_k)u^k\\
&=\sum_k\left(\sum_{i+j=k}\sigma(a_i)\sigma(b_j)\right)u^k\\
&=\left(\sum_i\sigma(a_i)u^i\right)
  \left(\sum_j\sigma(b_j)u^j\right)\\
&=\sigma_u(f)\sigma_u(g).
\end{align*}
所以$\sigma_u$为环同态。而且
\[
\sigma_u(x)=\sigma(1_R)u=u,\qquad
\sigma_u(a)=\sigma(a)\quad(a\in R),
\]
故它是所需扩充。

任一扩充同态在$R$上的取值由$\sigma$确定，在$x$上的取值又已指定为$u$，因而必须满足
\[
\sigma_u\left(\sum_i a_i x^i\right)=\sum_i\sigma(a_i)u^i.
\]
所以扩充唯一，其像恰为$\sigma(R)[u]$。
\end{proof}

\begin{corollary}
设$S$为$R$的扩环，$u\in S$。存在$R[x]$的理想$I$，使
\[
R[u]\cong R[x]/I,\qquad I\cap R=(0).
\]
同构由$f(x)+I\mapsto f(u)$给出。
\end{corollary}
\begin{proof}
在上一定理中取$\sigma:R\hookrightarrow S$为包含映射。其扩充$\sigma_u:f(x)\mapsto f(u)$的像为$R[u]$。令$I=\ker\sigma_u$，由环同态基本定理，$R[x]/I\cong R[u]$。

若$a\in I\cap R$，则一方面$\sigma_u(a)=0$，另一方面$\sigma_u(a)=a$，所以$a=0$。故$I\cap R=(0)$。
\end{proof}

\begin{definition}
映射$f(x)\mapsto f(u)$称为\textbf{代入同态}。若$f(u)=0$，称$u$为$f(x)$的\textbf{根}，并称$f(u)=0$为$u$在$R$上的一个代数关系。

若不存在非零$f\in R[x]$使$f(u)=0$，称$u$在$R$上\textbf{超越}；若存在这样的非零多项式，则称$u$在$R$上\textbf{代数}，或称$u$为$R$上的代数元。
\end{definition}

因此$I$恰记录$u$的全部代数关系。$I=(0)$时，$R[u]\cong R[x]$；$I\ne(0)$时，$u$满足非平凡的代数关系。

\begin{proposition}
若$I$为$R[x]$的理想且$I\cap R=(0)$，则自然同态$R[x]\to R[x]/I$限制在$R$上为单同态。把$R$嵌入此商环，并记$u=x+I$，便有$R[x]/I=R[u]$。
\end{proposition}
\begin{proof}
限制映射$a\mapsto a+I$的核为$I\cap R=(0)$，故为单射，并保持单位元。商环中任一元素都可写成
\[
f(x)+I=(a_0+I)+(a_1+I)u+\cdots+(a_n+I)u^n.
\]
将$a_i+I$与$a_i$等同后，这就是$R[u]$中的元素，故商环由$R$与$u$生成。
\end{proof}

\begin{definition}
对正整数$n$归纳定义\textbf{多元多项式环}：$n=1$时即$R[x_1]$；$n>1$时，
\[
R[x_1,\ldots,x_n]=R[x_1,\ldots,x_{n-1}][x_n].
\]
它的元素可写成有限和
\[
f(x_1,\ldots,x_n)
=\sum_{i_1,\ldots,i_n\ge0}a_{i_1\ldots i_n}x_1^{i_1}\cdots x_n^{i_n},
\qquad a_{i_1\ldots i_n}\in R.
\]
其中$x_1^{i_1}\cdots x_n^{i_n}$称为单项式。
\end{definition}

上述表示的唯一性由归纳得到。先按$x_n$展开，$f=\sum_j f_j(x_1,\ldots,x_{n-1})x_n^j$。由一元情形，$f=0$当且仅当所有$f_j=0$；再由归纳假设，每个$f_j=0$当且仅当它的各个系数均为零。因此
\[
f(x_1,\ldots,x_n)=0
\quad\Longleftrightarrow\quad
a_{i_1\ldots i_n}=0\text{ 对所有指数组成立}.
\]

\begin{theorem}\label{p2:multi-poly-universal}
设$R,S$为交换幺环，$\sigma:R\to S$为环同态，且$\sigma(1_R)=1_S$。给定$u_1,\ldots,u_n\in S$，$\sigma$唯一地扩充为环同态
\[
\sigma_n:R[x_1,\ldots,x_n]\to S,
\qquad \sigma_n(x_i)=u_i\quad(1\le i\le n).
\]
具体地，
\[
\sigma_n(f)=\sum_{i_1,\ldots,i_n}
\sigma(a_{i_1\ldots i_n})u_1^{i_1}\cdots u_n^{i_n}.
\]
\end{theorem}
\begin{proof}
对$n$归纳。$n=1$为定理\ref{p2:poly-universal}。设结论对$n-1$成立，则已有唯一扩充
\[
\sigma_{n-1}:R[x_1,\ldots,x_{n-1}]\to S,
\qquad x_i\mapsto u_i\quad(1\le i<n).
\]
以$R[x_1,\ldots,x_{n-1}]$为系数环，再应用一元情形，指定$x_n\mapsto u_n$，得到所需的$\sigma_n$。逐次展开便得所列公式。

若另有同态满足条件，它在前$n-1$个变量的多项式环上的限制必为$\sigma_{n-1}$，再由一元扩充的唯一性，它在整个环上必与$\sigma_n$相同。因此扩充唯一。
\end{proof}

\begin{definition}
设$R'$为$R$的扩环，$u_1,\ldots,u_n\in R'$。集合
\[
R[u_1,\ldots,u_n]
=\{f(u_1,\ldots,u_n):f\in R[x_1,\ldots,x_n]\}
\]
称为这些元素在$R$上生成的子环，也称$R$上的\textbf{有限生成环}。

非零多项式$f\in R[x_1,\ldots,x_n]$给出的等式$f(u_1,\ldots,u_n)=0$称为一个非平凡代数关系。若不存在这样的关系，称$u_1,\ldots,u_n$在$R$上\textbf{代数无关}；否则称它们\textbf{代数相关}。
\end{definition}

\begin{corollary}
对$R'$中的任意$u_1,\ldots,u_n$，代入同态的核
\[
I=\{f\in R[x_1,\ldots,x_n]:f(u_1,\ldots,u_n)=0\}
\]
满足
\[
R[x_1,\ldots,x_n]/I\cong R[u_1,\ldots,u_n],\qquad I\cap R=(0).
\]
进一步，$I=(0)$当且仅当$u_1,\ldots,u_n$代数无关。
\end{corollary}
\begin{proof}
将包含映射$R\hookrightarrow R'$扩充，所得代入同态的像恰为$R[u_1,\ldots,u_n]$，由同构定理得到商环同构。若$a\in I\cap R$，则$a$既映为$a$又映为$0$，所以$a=0$。最后，$I=(0)$恰好表示没有非零多项式给出零关系。
\end{proof}

\begin{proposition}
设$N$为$R$的理想，$\overline R=R/N$，用$N[x_1,\ldots,x_n]$表示所有系数均在$N$中的多项式。则它是理想，并有
\[
R[x_1,\ldots,x_n]/N[x_1,\ldots,x_n]
\cong\overline R[y_1,\ldots,y_n].
\]
\end{proposition}
\begin{proof}
将自然同态$\sigma:R\to\overline R$扩充为$\sigma_n$，并指定$x_i\mapsto y_i$。它逐项地将系数$a_{i_1\ldots i_n}$映为$a_{i_1\ldots i_n}+N$，所以每个目标多项式均有原像，$\sigma_n$满射。

由系数表达式的唯一性，$\sigma_n(f)=0$当且仅当$f$的每个系数均属于$N$，即
\[
\ker\sigma_n=N[x_1,\ldots,x_n].
\]
核为理想，环同态基本定理给出所述同构。
\end{proof}

有限生成环$R[u_1,\ldots,u_n]$的结构由关系理想$I$决定。若$I=(f_1,\ldots,f_r)$，则$I$中的每个关系都是$\sum_j h_jf_j$，其中$h_j\in R[x_1,\ldots,x_n]$。能否用有限个基本关系得到全部关系，取决于这个理想是否有限生成；这里不对任意系数环作此断言。
